% !TeX root = ../复分析笔记.tex
% 由主文档通过 \input 载入；请编译主文档。
\setcounter{chapter}{1} % 保留原书章号。
\chapter{习题解答与关键步骤}
\label{chap:ca-solutions}

本附录与各章习题逐项对应。计算题给出结果和主要推导，证明题指出假设的使用位置。建议先独立完成，再对照思路；得到同一个数值并不能代替对路径方向、收敛区域和定理条件的核查。

\section{复数与复平面}
对应第~\ref{sec:ca-plane-ex} 节。
\begin{enumerate}
 \item 乘以分母的共轭可得 $(1+2i)/(2-i)=i$，模为 $1$；原分子、分母的模均为 $\sqrt5$，所以模的商也为 $1$。
 \item $8i=8e^{i\pi/2}$，三个根为 $2e^{i(\pi/6+2k\pi/3)}$，$k=0,1,2$。相邻辐角相差 $2\pi/3$，三根位于半径为 $2$ 的圆上的正三角形顶点。
 \item 平方并整理得 $3x^2+10x+3y^2+3=0$，即 $(x+5/3)^2+y^2=16/9$。圆心为 $-5/3$，半径为 $4/3$。
 \item 设 $q=(z-w)/(z+w)$。由单位圆条件有 $\overline z=1/z$、$\overline w=1/w$，故 $\overline q=(w-z)/(w+z)=-q$，这等价于 $\operatorname{Re}q=0$。分母非零条件保证整个表达式有定义。
 \item 右半平面是凸区域，所以单连通。去心单位圆盘不是单连通，绕原点一周的圆具有非零绕数，在避开原点的连续收缩中绕数不变。裂平面单连通：任意点唯一写成 $re^{i\theta}$，$r>0$、$-\pi<\theta<\pi$，可通过 $H(s,re^{i\theta})=((1-s)r+s)e^{i(1-s)\theta}$ 连续收缩到 $1$，始终不碰支割线。
\end{enumerate}

以下继续给出本章其余习题的解答。
\begin{description}
 \item[习题~\ref{ex:ca-roots-sum}]
 $-4=4e^{i\pi}$，故四个根为 $\sqrt2e^{i(\pi/4+k\pi/2)}$（$0\le k<4$），即 $1+i,-1+i,-1-i,1-i$。若 $n$ 整除 $m$，各项 $\omega^{mk}$ 都为 $1$，和为 $n$；否则 $\omega^m\ne1$，等比求和给出 $(1-\omega^{mn})/(1-\omega^m)=0$。负整数 $m$ 也适用。
 \item[习题~\ref{ex:ca-parallelogram}]
 分别展开 $|z\pm w|^2=|z|^2+|w|^2\pm2\operatorname{Re}(z\overline w)$，相加即得结论。这是平行四边形法则：两条对角线长度的平方和等于四条边长度的平方和。
 \item[习题~\ref{ex:ca-domains}]
 单位圆盘和右半平面都是凸区域，因而连通且单连通。去心单位圆盘连通，但绕原点的圆周不能收缩，故不单连通。$\mathbb{C}\setminus\mathbb{R}$ 分为上、下半平面，不连通，因此不是单连通区域；两个连通分支各自单连通。
 \item[习题~\ref{ex:ca-separation}]
 对每个 $a\in K$ 选 $r_a>0$，使 $B(a,2r_a)\subset U$。取 $B(a_j,r_{a_j})$ 的有限子覆盖，令 $\delta=\min_jr_{a_j}>0$。若 $z\in K$，选择覆盖它的第 $j$ 个圆盘；对 $w\notin U$，有 $|w-a_j|\ge2r_{a_j}$，故 $|w-z|\ge|w-a_j|-|z-a_j|>r_{a_j}\ge\delta$。取下确界得到所求。
\end{description}



\section{复导数与全纯函数}
对应第~\ref{sec:ca-derivative-ex} 节。
\begin{enumerate}
 \item 定义域为 $\mathbb C\setminus\{0\}$，在其上全纯，$f'(z)=3z^2-2/z^2$。原点处函数无有限延拓，不能列入全纯区域。
 \item $\operatorname{Re}z$ 的分量满足 $u_x=1$、$v_y=0$，C--R 方程处处失败，因此没有复可导点。$z\overline z=x^2+y^2$ 的 C--R 方程仅在原点成立，而分量实可微，故只在原点复可导，导数为零。两函数作为从 $\mathbb R^2$ 到 $\mathbb R^2$ 的映射都处处光滑。
 \item 在完全包含于 $D$ 的小圆盘中，任意线段参数化为 $z(t)=a+t(b-a)$。链式法则给出 $(f\circ z)'(t)=f'(z(t))(b-a)=0$，实部、虚部的一元导数均为零，故端点值相同。因此 $f$ 局部为常数；任意两点可由区域内折线连接，逐段使用该结论得全局为常数。
 \item 写 $f=u+i0$，C--R 方程给出 $u_x=u_y=0$，沿区域内折线可知 $u$ 恒定。也可用开映射定理反证，但此处的证明不依赖后面的定理。
 \item 由 $v_y=u_x=3x^2-3y^2$，积分得 $v=3x^2y-y^3+\phi(x)$；再用 $v_x=-u_y=6xy$ 得 $\phi'=0$。所以 $v=3x^2y-y^3+C$，对应 $f=z^3+iC$。
\end{enumerate}

以下继续给出本章其余习题的解答。
\begin{description}
 \item[习题~\ref{ex:ca-affine}]
 任一点的差商都为 $a+b\overline h/h$。沿实轴与虚轴分别等于 $a+b$、$a-b$，故复可微当且仅当 $b=0$。此时处处可微且 $f'=a$；否则没有复可微点。
 \item[习题~\ref{ex:ca-c1-criterion}]
 连续偏导数满足柯西--黎曼方程的条件为 $2x=\beta x$、$2\alpha y=-\beta y$，须对所有 $x,y$ 成立。因此 $\beta=2$、$\alpha=-1$，此时 $f=z^2$、$f'=2z$。
 \item[习题~\ref{ex:ca-real-constant}]
 不必。定义域有两个连通分支，实值全纯函数在每个分支上各为常数，但两个常数可以不同。例如右分支取 $1$、左分支取 $0$，在每一点附近均为常数，因而全纯。
 \item[习题~\ref{ex:ca-modulus-constant}]
 $c=0$ 时 $f\equiv0$。若 $c>0$，写 $f=u+iv$，对 $u^2+v^2=c^2$ 分别求偏导，再用柯西--黎曼方程，得到 $u u_x+v v_x=0$、$v u_x-u v_x=0$。关于 $u_x,v_x$ 的系数行列式为 $-(u^2+v^2)=-c^2\ne0$，故二者均为零，即 $f'=0$。由命题~\ref{prop:ca-zero-derivative}，$f$ 为常数。
\end{description}



\section{初等函数与分支}
对应第~\ref{sec:ca-elementary-ex} 节。
\begin{enumerate}
 \item $1+i=\sqrt2e^{i\pi/4}$，故 $z=\tfrac12\ln2+i(\pi/4+2k\pi)$，$k\in\mathbb Z$。
 \item $\operatorname{Log}(1-i)=\tfrac12\ln2-i\pi/4$。$(-1)^{1/3}$ 的三个值为 $e^{i\pi/3},e^{i\pi},e^{5i\pi/3}$，即 $\tfrac12+i\tfrac{\sqrt3}2,-1,\tfrac12-i\tfrac{\sqrt3}2$。三次方程的全部解与任取一个对数分支后的函数值不是同一问题。
 \item 由 $(w-w^{-1})/(2i)=2$ 得 $w^2-4iw-1=0$，故 $w=i(2\pm\sqrt3)$。解 $e^{iz}=w$ 得
 \[
 z=\frac\pi2+2k\pi\pm i\ln(2+\sqrt3),\qquad k\in\mathbb Z.
 \]
 两种符号对应互为倒数的正数 $2\pm\sqrt3$，应保留全部整数周期。
 \item 右半平面上的辐角取在 $(-\pi/2,\pi/2)$。若 $z=re^{i\theta}$，则 $1/z=r^{-1}e^{-i\theta}$ 仍在右半平面，所以对数为 $-\ln r-i\theta$，正好等于 $-\operatorname{Log}z$。结论依赖这两个点使用相容分支。
 \item 假设 $s$ 存在，令 $q(t)=s(e^{it})e^{-it/2}$，则 $q(t)^2=1$。$q$ 连续且只取 $\{1,-1\}$，所以为常数。然而 $q(2\pi)=-s(1)=-q(0)$，而 $s(1)^2=1$ 保证 $q(0)\ne0$，矛盾。
\end{enumerate}

以下继续给出本章其余习题的解答。
\begin{description}
 \item[习题~\ref{ex:ca-exp-solutions}]
 $e^z=-2$ 的解为 $\log2+(2k+1)\pi i$。$\sin z=0$ 等价于 $e^{2iz}=1$，故 $z=k\pi$；$\cosh z=0$ 等价于 $e^{2z}=-1$，故 $z=(k+1/2)\pi i$，均取 $k\in\mathbb{Z}$。
 \item[习题~\ref{ex:ca-log-product}]
 记 $s=\operatorname{Arg}z+\operatorname{Arg}w\in(-2\pi,2\pi)$。若 $-\pi<s<\pi$，两式相等；若 $\pi<s<2\pi$，则 $\operatorname{Log}(zw)=\operatorname{Log}z+\operatorname{Log}w-2\pi i$；若 $-2\pi<s<-\pi$，则应加上 $2\pi i$。$s=\pm\pi$ 会使 $zw$ 位于负实轴，与题设不符。
 \item[习题~\ref{ex:ca-square-root-domain}]
 对 $z=re^{i\theta}$，$r>0$、$0<\theta<\pi$，取 $S(z)=\sqrt r\,e^{i\theta/2}$。其像恰为第一象限 $\{w:\operatorname{Re}w>0,\operatorname{Im}w>0\}$，并有 $S'(z)=1/(2S(z))$。像集的满性由将第一象限中的数平方得到。
\end{description}



\section{复曲线积分}
对应第~\ref{sec:ca-integral-ex} 节。
\begin{enumerate}
 \item 取全平面上的原函数 $z^3/3$，结果为 $(1+i)^3/3=(-2+2i)/3$，与路径无关。
 \item 参数化 $z=e^{it}$ 后有 $\overline z\,\mathrm dz=i\,\mathrm dt$，积分为 $2\pi i$。在单位圆上 $\overline z=1/z$，由基本圆周积分得到同样结果。这个相等只在单位圆上成立，不是两个函数在某个区域内恒等。
 \item 上半圆取 $z=e^{it}$、$t:0\to\pi$，结果为 $i\pi$；下半圆从 $1$ 到 $-1$ 可取 $t:0\to-\pi$，结果为 $-i\pi$。两者之差 $2\pi i$ 反映合成闭路绕原点一周。
 \item 在 $|z|=R$ 上，反三角不等式给出 $|z^2+1|\ge R^2-1$，半圆长度为 $\pi R$，代入长乘高估计得所需上界，并随 $R\to\infty$ 趋于零。
 \item $(F-G)'=0$，由区域连通和前章习题结论得差为常数。若定义域不连通，可以在不同连通分支取不同常数，不能保证全域只有一个常数。
\end{enumerate}

以下继续给出本章其余习题的解答。
\begin{description}
 \item[习题~\ref{ex:ca-circle-conjugate}]
 圆周上 $\overline z=R^2/z$，故积分为 $R^2\int_\gamma dz/z=2\pi iR^2$；反向时为 $-2\pi iR^2$。也可直接用 $z=Re^{it}$ 参数化检验。
 \item[习题~\ref{ex:ca-log-integral}]
 右半平面上的 $1/z$ 有原函数 $\operatorname{Log}z$，故积分为 $\operatorname{Log}(1+i)-\operatorname{Log}1=\tfrac12\log2+\tfrac{\pi i}{4}$，与路径无关。
 \item[习题~\ref{ex:ca-winding-path}]
 差商 $\gamma'(t)/(\gamma(t)-a)=ik$，积分为 $2\pi ik$，所以绕数是 $k$。$k=0$ 时曲线恒等于 $a+r$，导数为零，绕数也是零。
 \item[习题~\ref{ex:ca-arc-estimate}]
 $ML$ 估计给出积分模不超过 $\pi C R^{1-p}$，当 $p>1$ 时趋于零。对上半圆 $z=Re^{it}$、$0\le t\le\pi$，$\int dz/z=\int_0^\pi i\,dt=\pi i$，说明 $p=1$ 时仅有这一量级控制不足。
\end{description}



\section{柯西积分理论}
对应第~\ref{sec:ca-cauchy-ex} 节。
\begin{enumerate}
 \item 分母为 $z^2$，对应 $\sin z$ 在零点的一阶导数，所以积分为 $2\pi i\cos0=2\pi i$。
 \item 部分分式为 $-1/(z-1)+2/(z-2)$。两个点均在圆内，故积分为 $2\pi i(-1+2)=2\pi i$。
 \item 在原点用半径 $R$ 的柯西估计，得 $|f^{(n)}(0)|\le n!C(1+R^m)/R^n$。对 $n>m$ 令 $R\to\infty$，得 $f^{(n)}(0)=0$。由整函数的泰勒展开，$f$ 是次数不超过 $m$ 的多项式。这里用到了第~\ref{chap:ca-series} 章证明的泰勒定理；读到本题时也可先保留为后续应用。
 \item 若 $\operatorname{Re}f\le A$，则 $|e^f|\le e^A$，所以 $e^f$ 由刘维尔定理为常数。微分得 $e^ff'=0$，又 $e^f\ne0$，故 $f'=0$，于是 $f$ 为常数。
 \item 取圆环内任意正向圆 $|z|=r$，$1<r<2$，有 $\oint\mathrm dz/z=2\pi i\ne0$。若存在原函数，该积分必须为零，矛盾。圆环中的孔洞使全纯性不足以保证原函数存在。
\end{enumerate}

以下继续给出本章其余习题的解答。
\begin{description}
 \item[习题~\ref{ex:ca-cauchy-integrals}]
 由高阶导数公式，第一项为 $(2\pi i/2!)(e^z)''|_{z=0}=\pi i$，第二项为 $2\pi i(\cos z)'|_{z=1}=-2\pi i\sin1$。两处分母零点都严格位于路径内部。
 \item[习题~\ref{ex:ca-derivative-bound}]
 对任意 $0<r<R-|a|$，闭圆盘 $\overline{B(a,r)}$ 包含于 $B(0,R)$，故 $|f'(a)|\le M/r$。令 $r\uparrow R-|a|$ 得到结论；无需把积分路径放到原定义域的边界上。
 \item[习题~\ref{ex:ca-entire-bounded}]
 若全平面上 $|f|\le M$，则对任意 $a$、任意 $R>0$ 都有 $|f'(a)|\le M/R$。令 $R\to\infty$ 得 $f'(a)=0$，故 $f$ 为常数。在固定圆盘中不能令可用半径趋于无穷大；例如 $f(z)=z$ 在单位圆盘内有界而非常数。
 \item[习题~\ref{ex:ca-uniform-holomorphic}]
 每点有一个闭圆盘包含于 $U$，在该圆盘上一致收敛使极限连续，故 $f$ 在 $U$ 连续。对任意闭三角形 $T\subset U$，在其边界上一致收敛允许交换极限与积分，于是 $\int_{\partial T}f\,dz=\lim_n\int_{\partial T}f_n\,dz=0$。莫雷拉定理遂给出全纯性；若不先验证连续性，便缺少该定理的一项假设。
\end{description}



\section{幂级数与全纯函数结构}
对应第~\ref{sec:ca-series-ex} 节。
\begin{enumerate}
 \item 根值或比值公式给出 $R=1$。从 $\sum_{n=0}^\infty z^n=1/(1-z)$ 在圆盘内逐项微分并乘以 $z$，得到 $\sum_{n=1}^\infty nz^n=z/(1-z)^2$。
 \item 写 $2-z=1-(z-1)$，得 $1/(2-z)=\sum_{n=0}^\infty(z-1)^n$，半径为 $1$，最近奇点为 $z=2$。
 \item $\sin z-z=-z^3/6+O(z^5)$，故为三阶零点。$e^z-1=z+O(z^2)$，平方后首项为 $z^2$，故为二阶零点。
 \item 第一列零点在内部点 $0$ 聚集，恒等定理给出 $f\equiv0$。第二列点聚于边界点 $1$，不满足定理。事实上 $\sin(\pi/(1-z))$ 在单位圆盘全纯、并非恒零，且在每个 $1-1/n$ 处为零。
 \item 施瓦茨引理的导数等号情形给出 $f(z)=e^{i\theta}z$；由 $f'(0)=1$ 得 $e^{i\theta}=1$，故 $f(z)=z$。
 \item $|z^2+1|\le|z|^2+1\le2$。要取等号，须 $|z|=1$ 且 $z^2$ 与 $1$ 同向，故 $z^2=1$。最大值为 $2$，恰在 $z=1,-1$ 取得。
\end{enumerate}

\section{洛朗级数与奇点}
对应第~\ref{sec:ca-laurent-ex} 节。
\begin{enumerate}
 \item 在 $0<|z|<1$，
 \[
 \frac1{z(z-1)}=-\frac1z\frac1{1-z}=-\sum_{n=0}^\infty z^{n-1}.
 \]
 在 $|z|>1$，
 \[
 \frac1{z(z-1)}=\frac1{z^2}\frac1{1-z^{-1}}
 =\sum_{n=0}^\infty z^{-n-2}.
 \]
 两个展开的 $z^{-1}$ 系数分别为 $-1$、$0$，对应积分圆围住不同的奇点集合。
 \item $(1-\cos z)/z^2=1/2-z^2/24+\cdots$，原点可去，延拓值 $1/2$。$\sin z/z^3=z^{-2}-1/6+\cdots$，原点为二阶极点。$\sin(1/z)=z^{-1}-z^{-3}/3!+\cdots$，有无穷多负次幂，原点为本性奇点。
 \item 令 $g(z)=(z-a)f(z)$。由假设，$g$ 在 $a$ 可去，延拓值为零。其泰勒展开可提出因子 $z-a$，即 $g(z)=(z-a)h(z)$，其中 $h$ 全纯；在去心邻域内 $f=h$，因此 $f$ 也可去。这说明增长严格慢于 $1/|z-a|$ 的孤立奇点不能留下任何主部。
 \item $z^3+1$ 在无穷远点有三阶极点；$z/(z^2+1)$ 在无穷远点可去且值为零，因为换元后为 $\zeta/(1+\zeta^2)$；$\sin z$ 换元后为 $\sin(1/\zeta)$，故无穷远点为本性奇点。
 \item 对数的一个分支没有包含原点周围整圈的去心圆盘作为定义域，全局连续分支也不存在，所以不满足单值孤立奇点分类的前提。$e^{1/z}$ 在完整去心圆盘上单值全纯，才可用洛朗主部分类为本性奇点。
\end{enumerate}

\section{留数与实积分}
对应第~\ref{sec:ca-residue-ex} 节。
\begin{enumerate}
 \item 从 $\cos z/z^3=z^{-3}-\tfrac12z^{-1}+\cdots$ 得原点留数 $-1/2$。对第二个函数消去二阶因子，求导得
 \[
 \operatorname{Res}\left(\frac z{(z-1)^2(z+1)},1\right)
 =\left.\frac{\mathrm d}{\mathrm dz}\frac z{z+1}\right|_{z=1}
 =\frac14.
 \]
 \item 在 $1,2$ 的留数分别为 $-1,4$，两个点都在圆内，故积分为 $2\pi i\cdot3=6\pi i$。也可先做多项式除法，但常数项的闭路积分为零。
 \item 上半平面只有 $i$ 这个二阶极点，其留数为
 \[
 \left.\frac{\mathrm d}{\mathrm dz}(z+i)^{-2}\right|_{z=i}
 =-\frac2{(2i)^3}=\frac1{4i}.
 \]
 圆弧上函数为 $O(R^{-4})$，因此弧积分趋零，实积分等于 $2\pi i/(4i)=\pi/2$。
 \item $a=3,b=2$ 满足 $a>|b|$，结果为 $2\pi/\sqrt{9-4}=2\pi/\sqrt5$。
 \item 全实轴积分为 $\pi e^{-ab}/a$，被积函数是偶函数，故半轴结果为 $\pi e^{-ab}/(2a)$。$a>0$ 保证分母无实零点，$b>0$ 确定向上闭合的方向。
 \item 令 $t=x^q$，则 $x^{p-1}\,\mathrm dx=(1/q)t^{p/q-1}\,\mathrm dt$。由于 $0<p/q<1$，由钥匙孔积分公式得
 \[
 \int_0^\infty\frac{x^{p-1}}{1+x^q}\,\mathrm dx
 =\frac\pi q\csc\frac{\pi p}q.
 \]
 这两个不等式也分别是原点与无穷端的收敛条件。
 \item $1/z^2$ 在原点留数为零，但有二阶极点。$\operatorname{PV}\int_{-1}^1\mathrm dx/x=0$，而两侧普通反常积分分别发散，故普通积分不存在。两者都说明一个有限的积分相关数值并不代表函数局部行为正常。
\end{enumerate}

\section{零点计数}
对应第~\ref{sec:ca-counting-ex} 节。
\begin{enumerate}
 \item 在 $|z|=1$ 上，$|z^4+1|\le2<3=|3z|$。以 $3z$ 为主项应用鲁歇定理，单位圆内按重数恰有一个零点。
 \item 在 $|z|=2$ 上，$|z+1|\le3<16=|z^4|$，故有四个零点。严格不等式还给出 $|z^4+z+1|\ge16-3>0$，所以边界没有根。
 \item 圆内有三重零点 $1$ 与二阶极点 $-1$，故积分为 $2\pi i(3-2)=2\pi i$。也可直接写 $f'/f=3/(z-1)-2/(z+1)$ 后积分。
 \item 由定理~\ref{thm:ca-holomorphic-log} 选全纯对数 $L$，令 $g=e^{L/2}$，则 $g^2=f$。若 $h^2=f$，因 $f$ 无零点，有 $(h/g)^2=1$；该商连续、区域连通，故 $h/g$ 恒为 $1$ 或 $-1$，全部平方根为 $g$ 与 $-g$。
 \item 开映射定理说明非常数全纯函数的像必须为非空开集；直线和圆周在平面上都不包含任何开圆盘，矛盾。
\end{enumerate}

\section{共形映射}
对应第~\ref{sec:ca-conformal-ex} 节。
\begin{enumerate}
 \item $T(z)=1/(1-z)$ 满足 $T(0)=1$、$T(1)=\infty$、$T(\infty)=0$。由三点确定性，它就是唯一答案。
 \item 可取 $T(z)=(z-2i)/(z+2i)$。模平方差为 $|z+2i|^2-|z-2i|^2=8\operatorname{Im}z$，所以映到单位圆盘，且 $T(2i)=0$；由莫比乌斯反函数可核实双射。
 \item 第一象限通过 $\zeta=z^2$ 双全纯映到上半平面，$1+i$ 映到 $2i$。再用上一题变换，得到 $T(z)=(z^2-2i)/(z^2+2i)$。平方函数在该象限单射，因为两个相反的非零点不可能同时位于第一象限。
 \item 若右半平面的两点满足 $z_1^2=z_2^2$，则 $z_1=z_2$ 或 $z_1=-z_2$；后者不可能，故单射。辐角从 $(-\pi/2,\pi/2)$ 加倍到 $(-\pi,\pi)$，像为 $\mathbb C\setminus(-\infty,0]$。整个去心平面同时含 $z$ 与 $-z$，因此不单射。
 \item 圆盘自同构固定零点，必为 $e^{i\theta}z$。导数 $e^{i\theta}$ 既为正实数又模为 $1$，故等于 $1$，映射为恒等映射。
\end{enumerate}

\section{调和函数}
对应第~\ref{sec:ca-harmonic-ex} 节。
\begin{enumerate}
 \item $\Delta(x^2+y^2)=4$，所以不调和。$\Delta(e^x\cos y)=e^x\cos y-e^x\cos y=0$，调和；一个调和共轭为 $e^x\sin y$，对应 $e^z$。
 \item 对 $x^3-3xy^2$ 的共轭为 $3x^2y-y^3$，对 $2y$ 的共轭为 $-2x$。所以总共轭为 $3x^2y-y^3-2x+C$，相应全纯函数为 $z^3-2iz+iC$，$C\in\mathbb R$。
 \item 解为 $u(re^{i\theta})=2+r^3\cos(3\theta)+2r\sin\theta$，也就是 $2+x^3-3xy^2+2y$，中心值为 $2$。中心值等于边界平均值，非零频率的平均均为零。
 \item 分别对 $u$ 和 $-u$ 应用边界最大值原理，得 $u\le M$ 与 $-u\le M$，合并得 $|u|\le M$。不能直接把实调和函数的模误认为另一个调和函数。
 \item 由 $u_x=x/(x^2+y^2)$、$u_y=y/(x^2+y^2)$，再求导得 $u_{xx}+u_{yy}=0$。候选共轭的微分为
 \[
 -\frac y{x^2+y^2}\,\mathrm dx+\frac x{x^2+y^2}\,\mathrm dy.
 \]
 沿 $x=\cos t,y=\sin t$ 积分等于 $\int_0^{2\pi}\mathrm dt=2\pi$，非零增量阻止它成为全局单值函数的全微分。
\end{enumerate}

